%% ===================================================================
%%  Curriculum Vitae -- Daniel Martín Guerrero
%%  BSc in Physics, Universidad de Salamanca
%%  Compile with: pdflatex Daniel_Martin_Guerrero_CV.tex  (run twice)
%% ===================================================================
\documentclass[10pt,a4paper]{article}

% ---------- Encoding, fonts and machine-readable text --------------
\usepackage[T1]{fontenc}
\usepackage[utf8]{inputenc}
\usepackage[english]{babel}
\usepackage[sups,scaled=0.98]{XCharter}
\usepackage[xcharter,vvarbb]{newtxmath}
\usepackage[scaled=0.86]{sourcesanspro}
\usepackage{microtype}
\input{glyphtounicode}
\pdfgentounicode=1

% ---------- Page layout --------------------------------------------
\usepackage[a4paper,top=1.45cm,bottom=1.45cm,left=1.65cm,right=1.65cm]{geometry}
\pagestyle{empty}
\setlength{\parindent}{0pt}
\setlength{\parskip}{0pt}
\raggedbottom

% ---------- Colours -------------------------------------------------
\usepackage[dvipsnames]{xcolor}
\definecolor{accent}{HTML}{1B3A5C}   % deep "Prussian" blue
\definecolor{accentlight}{HTML}{4F7CAC}
\definecolor{muted}{HTML}{5A6270}
\definecolor{rulegrey}{HTML}{C9D1DB}

% ---------- Packages -----------------------------------------------
\usepackage{fontawesome5}
\usepackage{titlesec}
\usepackage{enumitem}
\usepackage{tabularx}
\usepackage{array}
\usepackage{needspace}
\usepackage[hidelinks,pdfusetitle]{hyperref}
\hypersetup{
  pdftitle={Daniel Martín Guerrero -- Curriculum Vitae},
  pdfauthor={Daniel Martín Guerrero},
  pdfsubject={Curriculum Vitae -- BSc in Physics, Universidad de Salamanca},
  pdfkeywords={Physics, Optics, Photonics, Astrophysics, Science Communication, Universidad de Salamanca}
}

% ---------- Section headings ----------------------------------------
\titleformat{\section}
  {\sffamily\bfseries\large\color{accent}}
  {}{0pt}
  {\MakeUppercase}
  [\vspace{-0.55ex}{\color{rulegrey}\titlerule[0.8pt]}]
\titlespacing*{\section}{0pt}{1.35ex}{0.9ex}

% ---------- Lists ---------------------------------------------------
\setlist[itemize]{
  leftmargin=1.15em, labelsep=0.5em,
  itemsep=0.15ex, topsep=0.35ex, parsep=0pt, partopsep=0pt,
  label={\color{accentlight}\scriptsize$\blacktriangleright$}
}

% ---------- Entry macros --------------------------------------------
% \entry{Title}{Date}{Subtitle}{Location}
\newcommand{\entry}[4]{%
  \needspace{6\baselineskip}%
  {\bfseries #1}\hfill{\sffamily\small\color{accent}#2}\par
  \if\relax\detokenize{#3}\relax\else
    \vspace{0.1ex}{\itshape\color{muted}#3}\hfill{\sffamily\small\color{muted}#4}\par
  \fi
}
% \award{Text}{Date}
\newcommand{\award}[2]{%
  \item\begin{tabularx}{\dimexpr\linewidth}[t]{@{}X@{\hspace{1em}}r@{}}
    #1 & {\sffamily\small\color{accent}#2}
  \end{tabularx}%
}
% Skill / coursework rows: hanging label column
\newlength{\skilllabelwidth}\setlength{\skilllabelwidth}{3.2cm}
\newcommand{\skillrow}[2]{%
  \noindent\parbox[t]{\skilllabelwidth}{\raggedright\leavevmode\sffamily\bfseries\small\color{accent}\strut#1\strut}%
  \hspace{0.8em}%
  \parbox[t]{\dimexpr\linewidth-\skilllabelwidth-0.8em}{\strut#2\strut}\par\vspace{0.55ex}%
}
\newenvironment{skilltable}{\par\vspace{0.2ex}}{\par}
\newcommand{\tech}[1]{\par\vspace{0.2ex}{\small\color{muted}\textit{Tools:} #1}}
\newcommand{\entrygap}{\vspace{0.9ex}}

\newcommand{\contact}[3]{%
  \href{#2}{\makebox[1.2em][c]{\color{accent}#1}\,#3}%
}

% ===================================================================
\begin{document}

% ------------------------------ Header ------------------------------
\begin{minipage}[t]{0.60\linewidth}
  \vspace{0pt}
  {\fontsize{25}{28}\selectfont\scshape\color{accent} Daniel Martín Guerrero}\\[1.0ex]
  {\large BSc in Physics \textcolor{muted}{(2025--2029)}}\\[0.35ex]
  {\color{muted} Faculty of Sciences, Universidad de Salamanca}\\[0.15ex]
  {\color{muted} Salamanca, Spain}
\end{minipage}%
\hfill
\begin{minipage}[t]{0.37\linewidth}
  \vspace{0.35ex}
  \raggedleft\small\sffamily
  \contact{\faPhone*}{tel:+34644419365}{+34-644\,419\,365}\\[0.55ex]
  \contact{\faEnvelope[regular]}{mailto:daniel.martgue@usal.es}{daniel.martgue@usal.es}\\[0.55ex]
  \contact{\faTiktok}{https://www.tiktok.com/@_.danielmg}{@\_.danielmg}\\[0.55ex]
  \contact{\faFilm}{https://letterboxd.com/danielmtg/}{letterboxd.com/danielmtg}
\end{minipage}

\vspace{1.4ex}
{\color{accent}\rule{\linewidth}{1.2pt}}\\[-0.6ex]
{\centering\small\sffamily\color{muted}
  Physics \,$\cdot$\, Optics \& Photonics \,$\cdot$\, Astrophysics \,$\cdot$\,
  Quantum Mechanics \,$\cdot$\, Science Communication\par}
\vspace{-0.2ex}

% ----------------------------- Education ----------------------------
\section{Education}

\entry{Universidad de Salamanca (USAL)}{2025 -- Present}
      {BSc in Physics -- second year, expected June 2029}{Salamanca, Spain}
\begin{itemize}
  \item Awarded \textbf{Honours} (\textit{Matrícula de Honor}) in \textit{Physics Laboratory}.
  \item Holder of the \textbf{University Access Award} (\textit{Premio de Acceso a la Universidad}),
        a full first-year tuition waiver.
\end{itemize}
\entrygap

\entry{Upper Secondary Education (Bachillerato) -- Science and Technology track}{2023 -- 2025}
      {IES Federico García Bernalt}{Salamanca, Spain}
\begin{itemize}
  \item Graduated with \textbf{Honours} (\textit{Matrícula de Honor}); overall average \textbf{9.88/10}.
  \item University entrance examination (PAU): \textbf{13.378/14}.
\end{itemize}

% ------------------------- Relevant coursework ----------------------
\section{Relevant Coursework}
\begin{skilltable}
  \skillrow{Linear Algebra \& Geometry}{Matrix and tensor calculus; vector spaces; linear, bilinear and
    sesquilinear maps (endomorphisms, metrics, polarities, quadratic forms); dual spaces and tensors;
    affine and Euclidean geometry; foundations of differential geometry.}
  \skillrow{Mathematical Analysis}{Differential and integral calculus in one and several variables;
    vector calculus; sequences and series; Taylor expansions and analysis of functions; complex variable.}
  \skillrow{Differential Equations}{Modelling of physical and natural systems: Malthusian and logistic
    population dynamics, Gompertz learning curves, SIS epidemic models, Lotka--Volterra systems,
    RLC circuits via the Laplace transform and signal analysis via the Fourier transform;
    introduction to topology, stochastic processes and measure theory.}
  \skillrow{Physics}{Classical mechanics, oscillations and waves (fundamentals of optics),
    electromagnetism, thermodynamics and computational techniques in physics.}
  \skillrow{Experimental Physics}{Data treatment and uncertainty quantification in the laboratory and
    through simulation: statistical estimation, covariance and propagation of variances, weighted means,
    least-squares fitting with $\chi^2$ model assessment and coverage intervals (Welch--Satterthwaite).}
\end{skilltable}

% ----------------------- Projects & independent study ---------------
\section{Independent Study \& Projects}

\entry{Foundations of Quantum Mechanics and Quantum Computing}{2026 -- Present}
      {Independent study -- author of structured \LaTeX\ lecture notes}{}
\begin{itemize}
  \item Authored \textit{Lesson~1: Hilbert Spaces and the Dirac Formalism -- Mathematical Foundations and
        the Born Postulate}, covering sesquilinear forms and inner products, Cauchy completeness and
        Hilbert spaces, the space $L^2(\mathbb{R})$ and square integrability, the Cauchy--Schwarz inequality,
        bras and kets, Hermitian adjoints, observables and spectral decomposition, and the Born rule
        and interference of probability amplitudes.
  \item Studying the fundamentals of quantum computation: qubits, observables and quantum logic
        gates (e.g.\ Hadamard).
\end{itemize}
\entrygap

\entry{AI-Assisted Study, Simulation and Modelling Workflow}{2025 -- Present}
      {Personal project -- Mathematics \& Physics}{}
\begin{itemize}
  \item Use frontier and open-source language models to produce structured study notes and to build
        simulations and models of physical situations, from parametrised manifolds to linear
        least-squares fits, as well as to formalise and develop new ideas.
  \item Use NotebookLM as a learning aid for problem solving and for the clear visualisation of
        complex concepts.
\end{itemize}
\tech{Claude Opus 5.5, Gemini, open-source language models, NotebookLM, Python, \LaTeX.}

% ------------------- Science communication & outreach ---------------
\section{Science Communication \& Outreach}

\entry{\textit{Eureka} -- Science Outreach Programme, Radio USAL}{2026 -- Present}
      {Scriptwriter \& Presenter; creator of the segment \textit{Cosmosofía}}{Salamanca, Spain}
\begin{itemize}
  \item Write and present science-outreach content for the official radio station of the University of
        Salamanca, within the \textit{Radioacción Cósmica} slot of the Supernova astronomy association.
  \item Created \textit{Cosmosofía}, a segment bridging astrophysics and the philosophy of science.
        Episode~I, \textit{``A Sea of Four Hundred Billion Suns''}, explores our cosmic coordinates,
        stellar nucleosynthesis and Cecilia Payne's discovery, science as an open conversation, and
        Sagan's \textit{Pale Blue Dot}.
\end{itemize}
\entrygap

\entry{Physics \& Mathematics Communication on TikTok -- \href{https://www.tiktok.com/@_.danielmg}{@\_.danielmg}}{Ongoing}
      {Content creator}{}
\begin{itemize}
  \item Produce AI-assisted videos, from general outreach to specialised topics, ranging from
        astrophysics, electromagnetism and thermodynamics to Sylvester's law of inertia, linear
        differential forms and the generalised Stokes theorem.
\end{itemize}
\entrygap

\entry{European Researchers' Night 2026 -- UCC+i, Universidad de Salamanca}{25 Sep.\ 2026}
      {Outreach activity \textit{``Houston, tenemos un telescopio''} (\textit{Houston, we have a telescope})}{Salamanca, Spain}
\begin{itemize}
  \item Delivered a three-hour public astronomy activity within the Scientific Culture and Innovation
        Unit (UCC+i) programme 2025--2026 (certified).
\end{itemize}

% --------------------- Talks & public engagement --------------------
\section{Talks \& Public Engagement}

\entry{Invited Speaker \& Round-Table Panellist -- XXI Philosophical Olympiad of Castilla y León}{27 Feb.\ 2026}
      {Universidad de Salamanca -- theme: \textit{Madness: Don Quixote in Today's World}}{Salamanca, Spain}
\begin{itemize}
  \item Delivered the talk \textit{``El naufragio de la luz. Disertación filosófica acerca de la lucidez
        del delirio. Don Quijote, Myshkin y la estética de la bondad''} (\textit{The Shipwreck of Light:
        a philosophical dissertation on the lucidity of delirium -- Don Quixote, Myshkin and the
        aesthetics of goodness}); certified by the Philosophical Olympiad Association of Castilla y León.
\end{itemize}
\entrygap

\entry{Panellist -- Presentation of \textit{Salvar a Sócrates} (Eduardo Infante, Ariel)}{5 Jun.\ 2026}
      {Librería Víctor Jara}{Salamanca, Spain}
\begin{itemize}
  \item Public conversation with the philosopher and author Eduardo Infante on his novel for learning
        to philosophise, alongside fellow panellists.
\end{itemize}

% -------------------------- Honours & awards ------------------------
\section{Honours \& Awards}
\begin{itemize}
  \award{\textbf{University Access Award}, Universidad de Salamanca -- granted to the 100 highest scores
         in the general phase of the 2024/25 university entrance examination}{2025}
  \award{\textbf{PAU score of 13.378/14}, Spanish university entrance examination}{2025}
  \award{\textbf{Honours} (\textit{Matrícula de Honor}), Bachillerato -- average 9.88/10}{2025}
  \award{\textbf{Honours} (\textit{Matrícula de Honor}) in \textit{Physics Laboratory}, BSc in Physics, USAL}{2026}
  \award{\textbf{Finalist}, Philosophical Olympiad of Castilla y León (OFCyL), philosophical dissertation
         category (Bachillerato) -- two consecutive editions (XIX and XX)}{2024, 2025}
  \award{\textbf{Participant}, PLANCKS 2027 -- international theoretical physics competition}{Upcoming}
\end{itemize}

% ------------------------ Memberships -------------------------------
\section{Memberships \& Affiliations}
\begin{itemize}
  \award{\textbf{Optica} (Advancing Optics and Photonics Worldwide) -- Individual Member}{Through Sep.\ 2027}
  \award{\textbf{OSAL Student Chapter} -- student association for optics outreach, Universidad de Salamanca}{Current}
  \award{\textbf{Supernova} University Astronomy Association -- outreach events, observation nights
         and retreats; radio scriptwriter and presenter}{2025 -- Present}
  \award{\textbf{Delegación Charra}, Student Group (GdeE) of the Royal Spanish Physics Society (RSEF)}{Current}
\end{itemize}

% -------------------------- Technical skills ------------------------
\section{Technical Skills}
\begin{skilltable}
  \skillrow{Programming}{Python, MATLAB, Arduino (C/C++-based language built on the Wiring framework),
    HTML (basic).}
  \skillrow{Electronics \& Lab}{Breadboard circuit prototyping with Arduino and LabVIEW; hands-on
    experimental work, measurement and instrumentation.}
  \skillrow{Scientific Tools}{Mathematica, PHYWE measure, LabVIEW, Arduino IDE,
    SAOImage DS9 (astronomical image analysis), \LaTeX\ (advanced).}
  \skillrow{AI Tools}{Claude Opus 5.5, Gemini, open-source language models, NotebookLM.}
  \skillrow{Languages}{Spanish (native); English (advanced -- trained up to C1 level; certified through
    the PAU English examination).}
\end{skilltable}

% ----------------------------- Interests ----------------------------
\section{Interests}
\begin{skilltable}
  \skillrow{Physics}{Astrophysics; optics and photonics (fundamental laser science, integrated
    photonics); particle physics; quantum mechanics and quantum computing.}
  \skillrow{Philosophy \& Film}{Philosophy of science and aesthetics; philosophy through cinema --
    Tarkovsky, Lynch, Bergman, Kurosawa and Masaki Kobayashi. Film reviewer on Letterboxd
    (\href{https://letterboxd.com/danielmtg/}{\color{accent}danielmtg}).}
\end{skilltable}

\end{document}
